% ── Chapter 5: Conclusion and Future Work ─────────────────────
\chapter{บทสรุปและงานในอนาคต (Conclusion and Future Work)}

% ─────────────────────────────────────────────────────────────
\section{บทสรุป (Conclusion)}
\label{sec:conclusion}
% ─────────────────────────────────────────────────────────────

\subsection{สรุปโปรเจกต์ (Summary of the Project)}

โปรเจกต์นี้ได้ออกแบบ, พัฒนา และประเมินผลเว็บแอปพลิเคชันที่ใช้ AI สำหรับการทดสอบความไวต่อยาปฏิชีวนะด้วยวิธี Kirby-Bauer disk diffusion (AST) แบบอัตโนมัติ ระบบนี้แก้ปัญหาคอขวดที่สำคัญในกระบวนการทำงานของห้องปฏิบัติการจุลชีววิทยาคลินิก นั่นคือการวัดขนาดบริเวณยับยั้งเชื้อ (ZOI) ด้วยไม้บรรทัดโดยเจ้าหน้าที่ ซึ่งใช้เวลานาน, ขึ้นอยู่กับดุลยพินิจของผู้ปฏิบัติงาน และยากต่อการตรวจสอบย้อนหลัง \cite{humphries2018} การทำให้ขั้นตอนนี้เป็นอัตโนมัติส่งผลโดยตรงต่อปริมาณงานที่ห้องปฏิบัติการรองรับได้, ความสามารถในการทำซ้ำของผลลัพธ์ และส่งผลดีต่อการรักษาผู้ป่วยในบริบทของวิกฤตการณ์การดื้อยาต้านจุลชีพทั่วโลก \cite{who_amr2022}

ระบบได้รวมองค์ประกอบทางเทคนิคหลักสามส่วนเข้าเป็น Pipeline แบบ End-to-end เพียงหนึ่งเดียว:

\begin{enumerate}
  \item \textbf{การตรวจจับวัตถุและการวัด ZOI (Object detection and ZOI measurement)} โดยใช้ YOLOv11n เป็นโมเดลหลัก \cite{wang2024yolov11} และ Faster R-CNN เป็นโมเดลรอง \cite{ren2015faster} พร้อมด้วยการใช้ Hough Transform แบบคลาสสิกเป็นเกณฑ์มาตรฐานที่ไม่ใช้การเรียนรู้ (No-learning baseline)
  \item \textbf{การระบุป้ายกำกับยาปฏิชีวนะ (Antibiotic label recognition)} โดยใช้ EasyOCR \cite{baek2019craft} พร้อมกับ Pipeline การประมวลผลเบื้องต้นแบบสองเส้นทาง (Otsu thresholding และ CLAHE)
  \item \textbf{การสอบเทียบพิกเซลเป็นมิลลิเมตรต่อรูปภาพ (Per-image pixel-to-mm calibration)} โดยใช้เส้นผ่านศูนย์กลางทางกายภาพขนาด 6\,มม. มาตรฐานของแผ่นยาปฏิชีวนะเป็นมาตราส่วนอ้างอิงภายในภาพ
\end{enumerate}

ผลลัพธ์ที่ประมวลผลแล้วจะถูกนำเสนอผ่านเว็บแอปพลิเคชันที่พัฒนาด้วย React ซึ่งแสดงขนาดเส้นผ่านศูนย์กลาง ZOI, การจำแนกประเภท S/I/R ตามเกณฑ์ CLSI และ EUCAST, ความสามารถในการแก้ไขผลด้วยตนเอง และการจัดเก็บประวัติผลการวิเคราะห์

\subsection{สรุปเชิงปริมาณ (Quantitative Summary)}

ผลลัพธ์เชิงปริมาณที่สำคัญของโปรเจกต์มีดังนี้

ในด้านความแม่นยำในการวัด ZOI โมเดล YOLOv11n แสดงให้เห็นถึงคุณภาพการตรวจจับที่เหนือกว่าด้วยค่า mAP@0.5 ที่ \textbf{95.91\%} และค่าความผิดพลาดสัมบูรณ์เฉลี่ย (MAE) ประมาณ \textbf{2.5--2.7\,มม.} ซึ่งมีประสิทธิภาพดีกว่า Faster R-CNN (MAE \textbf{4.65--4.85\,มม.}) และ Hough Transform อย่างมีนัยสำคัญ สถาปัตยกรรมแบบขั้นตอนเดียว (Single-stage) ของ YOLOv11n ให้ค่า Precision และ Recall ที่เกือบสมบูรณ์ ($\approx$1.0) พร้อมการลู่เข้าหาคำตอบที่รวดเร็ว

สำหรับการสอบเทียบ การนำ \emph{การสอบเทียบแบบไดนามิก (Dynamic Calibration)} มาใช้โดยใช้เส้นผ่านศูนย์กลางแผ่นยาปฏิชีวนะมาตรฐาน 6\,มม. เป็นตัวอ้างอิงในภาพ ประสบความสำเร็จในการลดความคลาดเคลื่อนเชิงระบบในการวัด วิธีนี้ช่วยให้แน่ใจว่าอัตราส่วนพิกเซลต่อมิลลิเมตรจะถูกคำนวณต่อแผ่นยาในทุกรูปภาพ ซึ่งสามารถปรับตามระยะห่างของกล้องถึงจานเพาะเชื้อและความบิดเบี้ยวของเลนส์ที่แตกต่างกันได้

สำหรับ OCR การใช้ EasyOCR ร่วมกับ Otsu thresholding สามารถบรรลุความแม่นยำประมาณ \textbf{95\%} ในการระบุป้ายกำหนดยาปฏิชีวนะ Pipeline การประมวลผลเบื้องต้นถูกระบุว่าเป็นปัจจัยหลักที่มีผลต่อประสิทธิภาพ โดยให้ผลลัพธ์ที่ดีกว่า TyphoonOCR ($\approx$85\%) และ EasyOCR แบบดิบ ($\approx$60\%)

สำหรับความเร็วในการประมวลผล Pipeline ที่รวมเข้าด้วยกันสามารถวิเคราะห์เสร็จสิ้นภายในเวลาประมาณ \textbf{3--5 วินาทีต่อรูปภาพ} ซึ่งรวมถึงขั้นตอนการตรวจจับ, OCR และการบันทึกข้อมูลลงฐานข้อมูล ซึ่งคิดเป็นการเพิ่มปริมาณงานได้มากกว่า 100 เท่าเมื่อเทียบกับกระบวนการวัดด้วยมือ

\subsection{สถานะความสำเร็จของโปรเจกต์ (Project Completion Status)}

ขณะนี้โปรเจกต์ \textbf{เสร็จสมบูรณ์ 100\%} และได้เปลี่ยนผ่านจากต้นแบบงานวิจัยไปสู่เครื่องมือทางคลินิกที่พร้อมใช้งานจริง Pipeline ทั้งหมด ซึ่งรวมถึงการจับภาพ, การตรวจจับ ZOI ด้วย YOLOv11n, การระบุป้ายกำกับด้วย EasyOCR และการสอบเทียบแบบไดนามิก ได้รับการรวมเข้าด้วยกันอย่างสมบูรณ์ ปัจจุบันระบบได้ \textbf{ติดตั้งใช้งานที่เซิร์ฟเวอร์ของโรงพยาบาลราชวิทยาลัยจุฬาภรณ์ (CRA)} โดยสามารถเข้าถึงได้ผ่านเครือข่ายทางคลินิกที่ปลอดภัยที่ \texttt{https://paniti-jupyter.cra.ac.th/zoneanalyzer/} การทดสอบการรวมระบบแบบ End-to-end เสร็จสิ้นแล้ว ซึ่งยืนยันว่า React frontend, FastAPI backend และฐานข้อมูล SQLite ทำงานร่วมกันได้อย่างราบรื่นในสภาพแวดล้อมการใช้งานจริง นอกจากนี้ระบบยังรองรับมาตรฐานการแปลผลทั้ง **CLSI 2025** และ **EUCAST 2023** โดยจะจับคู่ยาปฏิชีวนะที่ตรวจพบกับ Breakpoints ที่เกี่ยวข้องจากฐานข้อมูลโดยอัตโนมัติ

% ─────────────────────────────────────────────────────────────
\section{ข้อจำกัด (Limitations)}
\label{sec:limitations}
% ─────────────────────────────────────────────────────────────

การประเมินข้อจำกัดปัจจุบันอย่างตรงไปตรงมาเป็นสิ่งจำเป็นสำหรับการแปลผลลัพธ์ในบริบททางคลินิกและทางเทคนิคที่เหมาะสม

\subsection{ความคลาดเคลื่อนเชิงระบบที่เหลืออยู่ของการสอบเทียบ (Residual Calibration Bias)}

แม้ว่าการสอบเทียบแบบไดนามิกต่อแผ่นยาจะประสบความสำเร็จ แต่ความคลาดเคลื่อนเฉลี่ยประมาณ 1--2\,มม. ยังคงเกิดขึ้นได้ในสภาพแสงที่รุนแรงหรือรูปภาพที่มีมุมเอียงมาก สาเหตุหลักเกิดจากการใช้ Bounding box สี่เหลี่ยมที่ขนานกับแกน ซึ่งเป็นการประมาณค่าทางเรขาคณิตสำหรับบริเวณยับยั้งทรงกลม เมื่อจานเพาะเชื้อเอียงอย่างมีนัยสำคัญ บริเวณยับยั้งจะปรากฏเป็นรูปวงรี ทำให้ความกว้างของ Bounding box เบี่ยงเบนไปจากเส้นผ่านศูนย์กลางจริงเล็กน้อย

\subsection{ความหลากหลายของชุดข้อมูล (Dataset Diversity)}

โมเดลปัจจุบันได้รับการฝึกฝนด้วยชุดข้อมูลจำนวน 1,884 ภาพ (ภาพจริงและภาพสังเคราะห์) ซึ่งส่วนใหญ่เป็นเชื้อ \textit{S. aureus} และ \textit{E. coli} แม้ว่าจะมีความแม่นยำสูงมากสำหรับเชื้อชนิดเหล่านี้ แต่เชื้อชนิดอื่นที่มีความสำคัญทางคลินิกซึ่งมีรูปแบบการเจริญเติบโตเฉพาะตัว (เช่น \textit{Proteus} ที่มีการเจริญเติบโตแบบ Swarming หรือ \textit{Klebsiella} ที่มีความเหนียวมากหรือ Mucoid) อาจสร้างลักษณะบริเวณยับยั้งที่โมเดลยังไม่เคยพบในปริมาณมากพอ

% ─────────────────────────────────────────────────────────────
\section{งานในอนาคต (Future Work)}
\label{sec:future_work}
% ─────────────────────────────────────────────────────────────

ขอแนะนำการปรับปรุงต่อไปนี้สำหรับระยะต่อไปของการพัฒนา

\subsection{การศึกษาความถูกต้องทางคลินิกขนาดใหญ่ (Large-Scale Clinical Validation Study)}

ก่อนที่จะเสนอระบบเพื่อใช้ในทางคลินิกจริง จำเป็นต้องมีการศึกษาความถูกต้องแบบ Blinded validation study อย่างเป็นทางการ โดยเปรียบเทียบการวัด ZOI แบบอัตโนมัติกับมติของนักเทคนิคการแพทย์ที่มีประสบการณ์อย่างน้อยสามคนบนกลุ่มตัวอย่างเชื้อทางคลินิกที่หลากหลายที่ราชวิทยาลัยจุฬาภรณ์ สิ่งนี้จะให้หลักฐานทางสถิติเกี่ยวกับความสอดคล้องเชิงประเภท (Categorical agreement - CA) ของระบบกับวิธีแบบแมนนวลตามมาตรฐาน ISO 20776-2 \cite{iso20776_2} และเกณฑ์ประเมินในแนวทาง CLSI MM17 \cite{clsi_mm17} ซึ่งระบุอัตราที่ยอมรับได้ของความผิดพลาดร้ายแรงมาก (Very major errors), ความผิดพลาดร้ายแรง (Major errors) และความผิดพลาดเล็กน้อย (Minor errors) การประเมินเส้นทางการกำกับดูแล (เช่น FDA 510(k) หรือ CE-IVD) ควรได้รับการพิจารณาเป็นส่วนหนึ่งของการวางแผนนี้

\subsection{โมดูลระบาดวิทยาและแนวโน้มการดื้อยา (Epidemiology and Resistance Trend Module)}

ส่วนขยายในระยะยาวคือการเพิ่มโมดูลการวิเคราะห์ทางระบาดวิทยาที่รวบรวมผลลัพธ์การดื้อยาจากจานเพาะเชื้อทั้งหมดที่ประมวลผลและจัดเก็บไว้ในฐานข้อมูล อัตราการดื้อยาตามชนิดของเชื้อ, ช่วงเวลา และแผนกผู้ป่วยสามารถแสดงผลเป็น Dashboard แนวโน้ม และส่งออกเพื่อรายงานการควบคุมการติดเชื้อ สิ่งนี้จะเปลี่ยนระบบจากเครื่องมือวัดให้เป็นองค์ประกอบการจัดการข้อมูลในห้องปฏิบัติการที่มีคุณค่าต่อสาธารณสุข โดยตรงกับความต้องการการวิจัยลำดับความสำคัญของ WHO สำหรับการเฝ้าระวังการดื้อยาต้านจุลชีพ \cite{who_amr2022}
